% Apêndices, versão 8 (23 nov. 2026), a partir da v7. Os três instrumentos são PROPOSTAS, ainda não validadas nem aplicadas.
% V8: a configuração passou a incluir o modelo de linguagem, e o contexto documental foi simplificado, pois o Capítulo 5 ainda não define o nível de associação.
\begin{apendicesenv}

\partapendices

\chapter{Ficha de Registro de Auditoria (Proposta)}
\label{ap-ficha}

Esta ficha é uma proposta, a ser revista com a orientação antes da coleta. Ela será preenchida uma vez para cada execução. Os campos sem informação registram o motivo, que pode ser não coletado, indisponível, não se aplica ou falha.

\begin{table}[htbp]
\centering
\caption{Campos da ficha de registro de auditoria}
\label{tab-ficha}
\begingroup
\footnotesize
\SingleSpacing
\renewcommand{\arraystretch}{1.15}
\begin{tabularx}{\textwidth}{>{\raggedright\arraybackslash}p{3.4cm}>{\raggedright\arraybackslash}X}
\hline
\textbf{Bloco} & \textbf{Campos} \\ \hline
Identificação & Código da auditoria, código do repositório, \textit{commit}, critério de inclusão, data, código da execução e, se for reexecução, a execução de origem e o motivo \\ \hline
Configuração & Versões da solução, das ferramentas, do modelo de linguagem e do arquivo de regras, com a soma de verificação do arquivo de regras \\ \hline
Estado & Concluída, parcial ou com falha, com a descrição do ocorrido \\ \hline
Achados & Número de achados de qualidade, de segurança e de dependências, com a referência ao relatório original de cada ferramenta \\ \hline
Consultas P2P & Para cada consulta, o solicitante, o destinatário, o achado a que se refere, a situação da resposta e a conclusão antes e depois \\ \hline
Contexto documental & Para cada resposta do agente de Requisitos e Documentação, o documento, o trecho utilizado e o resultado da verificação do trecho pelo autor \\ \hline
Parecer & Categoria da recomendação, regras aplicadas e seu tipo, pontos de atenção, lacunas e questão encaminhada ao responsável humano \\ \hline
Correspondência & Para cada afirmação do parecer verificada pelo autor, a evidência citada e o resultado da verificação \\ \hline
Divulgação & Restrições de divulgação e local de guarda dos registros originais \\ \hline
\end{tabularx}
\endgroup
\legend{Fonte: elaboração do autor.}
\end{table}

\chapter{Instrumento de Apreciação dos Pareceres (Proposta)}
\label{ap-instrumento}

Este instrumento é uma proposta e não foi validado. Ele será respondido uma vez por especialista para cada parecer, depois da leitura do material descrito na seção~\ref{met-apreciacao}.

\section*{Parte 1. Perfil do participante}

Esta parte é respondida uma única vez.

\begin{enumerate}
    \item Área de atuação atual.
    \item Tempo de experiência em segurança de software, em anos.
    \item Relação com a Empresa A (nenhuma, colaborador ou prestador de serviço).
\end{enumerate}

\section*{Parte 2. Critérios}

Para cada afirmação, o participante indica seu grau de concordância em uma escala de cinco pontos, de discordo totalmente a concordo totalmente.

\begin{enumerate}
    \item Clareza. O parecer permite entender o que foi encontrado e o que está sendo recomendado.
    \item Coerência. A recomendação e os pontos de atenção decorrem das evidências apresentadas.
    \item Fundamentação. As regras e as fontes citadas sustentam o que o parecer afirma.
\end{enumerate}

\section*{Parte 3. Adequação da recomendação}

\begin{enumerate}
    \item Diante do material apresentado, a recomendação do parecer é adequada? As opções são sim, em parte, não e o material não é suficiente para julgar.
    \item Justifique a resposta. Se respondeu em parte ou não, indique o que mudaria na recomendação.
\end{enumerate}

\section*{Parte 4. Comentários}

\begin{enumerate}
    \item Alguma informação fez falta para avaliar o parecer?
    \item As consultas entre os agentes, mostradas no material, acrescentaram algo à sua leitura do achado?
    \item Outros comentários.
\end{enumerate}

\chapter{Roteiro do Termo de Consentimento Livre e Esclarecido (Proposta)}
\label{ap-tcle}

% PENDENTE: este roteiro lista o conteúdo mínimo. A redação final do TCLE e a necessidade de submissão a comitê de ética devem ser definidas com a orientação, conforme as normas da universidade.
Este roteiro é uma proposta do conteúdo do termo. A redação final e as exigências institucionais aplicáveis serão definidas com a orientação.

\begin{enumerate}
    \item Identificação da pesquisa, do pesquisador e dos orientadores, com os contatos.
    \item Objetivo da participação, que é apreciar pareceres produzidos por um projeto piloto quanto à clareza, à coerência e à fundamentação.
    \item O que será pedido ao participante, com a estimativa de tempo por parecer.
    \item Caráter voluntário da participação e possibilidade de desistência a qualquer momento, sem prejuízo.
    \item Registro das respostas sem identificação nominal e forma de guarda dos dados.
    \item Confidencialidade do material de leitura, que pode conter informações de repositórios da Empresa A.
    \item Uso dos dados, restrito à monografia e a publicações dela derivadas.
    \item Espaço para a assinatura e para a data.
\end{enumerate}

\end{apendicesenv}
